\begin{figure}[H]
    \centering
    \caption{O processo do Paradoxo de Banach-Tarski}
    \label{fig:banach_tarski_vertical}

    \begin{tikzpicture}[
            scale=0.9,
            transform shape,
            font=\sffamily,
            >=stealth
        ]

        % Cores Condicionais
        \def\ColorBall{\ifbool{darkmode}{gray!30}{gray!80}}
        \def\ColorA{\ifbool{darkmode}{DarkModeLink}{LightModeLink}}
        \def\ColorB{\ifbool{darkmode}{DarkModeViolet}{LightModeViolet}}
        \def\ColorText{\ifbool{darkmode}{DarkModeText}{black}}

        % ==========================================
        % ETAPA 1: Esfera Original
        % ==========================================
        \begin{scope}[shift={(0,0)}]
            \shade[ball color=\ColorBall, opacity=0.9] (0,0) circle (1.5);
            \draw[\ColorText, dashed] (1.5,0) arc (0:180:1.5 and 0.4);
            \draw[\ColorText] (1.5,0) arc (0:-180:1.5 and 0.4);
            \draw[\ColorText, thick] (0,0) circle (1.5);

            \node[\ColorText, font=\Large\bfseries] at (0,0) {\(S\)};
            \node[\ColorText, right=2cm, align=left] {
                \textbf{1. Esfera Original}\\
                Bola sólida em \(\mathbb{R}^3\)
            };
        \end{scope}

        \draw[->, \ColorText, ultra thick] (0, -1.8) -- (0, -2.8);

        % ==========================================
        % ETAPA 2: Esfera com cortes representativos
        % ==========================================
        \begin{scope}[shift={(0,-4.5)}]
            \shade[ball color=\ColorBall, opacity=0.9] (0,0) circle (1.5);

            % Linhas de corte simbólicas (representando a decomposição complexa)
            \draw[\ColorText, thick] (-1.2, 0.9) to[out=-45, in=135] (0.5, -1.4);
            \draw[\ColorText, thick] (-0.8, -1.2) to[out=60, in=200] (1.3, 0.7);
            \draw[\ColorText, thick] (0.2, 1.4) to[out=-90, in=180] (1.4, -0.5);
            \draw[\ColorText, thick, dashed] (-1.4, -0.5) to[out=30, in=150] (1.2, 0.9);

            \draw[\ColorText, thick] (0,0) circle (1.5);

            \node[\ColorText, right=2cm, align=left] {
                \textbf{2. Decomposição}\\
                Divisão em 5 pedaços\\
                (subconjuntos não-mensuráveis)
            };
        \end{scope}

        \draw[->, \ColorText, ultra thick] (0, -6.3) -- (0, -7.3);

        % ==========================================
        % ETAPA 3: Pedaços separados
        % ==========================================
        \begin{scope}[shift={(0,-9.5)}]
            % Subconjuntos tipo "A" (irão formar S1)
            \begin{scope}[shift={(-1.2, 0.8)}]
                \shade[ball color=\ColorA, opacity=0.8] plot [smooth cycle, tension=0.7] coordinates {(0,0) (0.8,0.2) (1,-0.5) (0.2,-0.8) (-0.4,-0.4)};
                \node[\ColorText] at (0.3, -0.2) {\(A_1\)};
            \end{scope}

            \begin{scope}[shift={(1.0, 0.5)}]
                \shade[ball color=\ColorA, opacity=0.8] plot [smooth cycle, tension=0.7] coordinates {(0,0) (0.5,0.6) (1.2,0) (0.6,-0.6)};
                \node[\ColorText] at (0.6, 0.1) {\(A_2\)};
            \end{scope}

            \begin{scope}[shift={(-0.8, -0.6)}]
                \shade[ball color=\ColorA, opacity=0.8] plot [smooth cycle, tension=0.7] coordinates {(0,0) (0.6,0.3) (0.8,-0.4) (-0.2,-0.7)};
                \node[\ColorText] at (0.4, -0.2) {\(A_3\)};
            \end{scope}

            % Subconjuntos tipo "B" (irão formar S2)
            \begin{scope}[shift={(1.2, -0.8)}]
                \shade[ball color=\ColorB, opacity=0.8] plot [smooth cycle, tension=0.7] coordinates {(0,0) (0.7,0.4) (1.0,-0.3) (0.4,-0.8)};
                \node[\ColorText] at (0.5, -0.1) {\(B_1\)};
            \end{scope}

            \begin{scope}[shift={(0, -1.2)}]
                \shade[ball color=\ColorB, opacity=0.8] plot [smooth cycle, tension=0.7] coordinates {(0,0) (0.9,0.1) (0.5,-0.6) (-0.3,-0.4)};
                \node[\ColorText] at (0.3, -0.2) {\(B_2\)};
            \end{scope}

            \node[\ColorText, right=2.8cm, align=left] {
                \textbf{3. Separação}\\
                As partes são desagrupadas
            };
        \end{scope}

        % Setas divergentes indicando os movimentos rígidos
        \draw[->, \ColorA, ultra thick, dashed] (-0.5, -11.5) to[out=-135, in=90]
        node[midway, left=0.5cm, align=right] {Isometrias\\(Rotação/Translação)} (-2.5, -13.5);

        \draw[->, \ColorB, ultra thick, dashed] (0.5, -11.5) to[out=-45, in=90]
        node[midway, right=0.5cm, align=left] {Isometrias\\(Rotação/Translação)} (2.5, -13.5);

        % ==========================================
        % ETAPA 4: Duas Cópias Originais
        % ==========================================
        % Cópia 1
        \begin{scope}[shift={(-2.5, -15.5)}]
            \shade[ball color=\ColorA, opacity=0.9] (0,0) circle (1.5);
            \draw[\ColorA, dashed] (1.5,0) arc (0:180:1.5 and 0.4);
            \draw[\ColorA] (1.5,0) arc (0:-180:1.5 and 0.4);
            \draw[\ColorA, thick] (0,0) circle (1.5);

            \node[\ColorText, font=\Large\bfseries] at (0,0) {\(S_1\)};
            \node[\ColorText, below=1.8cm, align=center] {Cópia 1 \\ \(S_1 \cong S\)};
        \end{scope}

        % Cópia 2
        \begin{scope}[shift={(2.5, -15.5)}]
            \shade[ball color=\ColorB, opacity=0.9] (0,0) circle (1.5);
            \draw[\ColorB, dashed] (1.5,0) arc (0:180:1.5 and 0.4);
            \draw[\ColorB] (1.5,0) arc (0:-180:1.5 and 0.4);
            \draw[\ColorB, thick] (0,0) circle (1.5);

            \node[\ColorText, font=\Large\bfseries] at (0,0) {\(S_2\)};
            \node[\ColorText, below=1.8cm, align=center] {Cópia 2 \\ \(S_2 \cong S\)};
        \end{scope}

        % Rótulo final centralizado
        \node[\ColorText, right=4.5cm, align=left] at (0, -15.5) {
            \textbf{4. Reconstrução}\\
            Duas esferas idênticas à original
        };

    \end{tikzpicture}

    \fonte{Elaborado pelo autor (2026).}
\end{figure}